% =============================================================
% บทที่ 4 - การวิเคราะห์ ออกแบบและพัฒนาระบบ
% โครงงานประเภท B (งานวิจัย)
% =============================================================

% -----------------------------------------------------------
\section{การวิเคราะห์และออกแบบสถาปัตยกรรม}
% -----------------------------------------------------------

\subsection{หลักการออกแบบเพื่อให้การเปรียบเทียบยุติธรรม}

คำถามวิจัยของโครงงานนี้คือความสามารถในการตัดสินใจเองของแบบจำลองภาษาให้ประโยชน์จริงหรือไม่ การจะตอบคำถามนี้ได้ การเปรียบเทียบต้องไม่มีตัวแปรอื่นปน หลักการออกแบบจึงกำหนดว่า\textbf{ทุกสถาปัตยกรรมต้องใช้ส่วนประกอบชุดเดียวกันทั้งหมด} และต่างกันเฉพาะชั้นที่ตัดสินใจเท่านั้น

หากสถาปัตยกรรมใดมีส่วนค้นคืนหรือส่วนจัดอันดับเป็นของตนเอง ผลที่วัดได้จะบอกไม่ได้ว่ามาจากสถาปัตยกรรมหรือมาจากส่วนประกอบที่เก่งกว่า การออกแบบนี้จึงเป็นเงื่อนไขที่ทำให้ข้อสรุปของโครงงานมีความหมาย

\subsection{ชั้นที่แบบจำลองภาษาตัดสินใจเอง}

ระบบที่พัฒนามีชั้นที่แบบจำลองภาษาตัดสินใจเองได้ 2 ชั้น และความต่างของทั้งสามสถาปัตยกรรมอยู่ที่ชั้นเหล่านี้เท่านั้น ดังตารางที่ \ref{tab:arms}

\begin{table}[H]
\centering
\caption{ชั้นที่แบบจำลองภาษาตัดสินใจเองในแต่ละสถาปัตยกรรม}
\label{tab:arms}
\begin{tabular}{|l|c|c|c|}
\hline
\textbf{ชั้นที่ตัดสินใจ} & \textbf{A1} & \textbf{A3C} & \textbf{A3} \\
\hline
เลือกว่าจะถามกลับ สนทนา หรือค้นหา & -- & \cmark & \cmark \\ \hline
ตัดสินว่าจะค้นซ้ำอีกรอบหรือไม่ & -- & กฎในโปรแกรม & \cmark \\ \hline
\textbf{จัดเป็นสถาปัตยกรรมแบบ} & \textbf{RAG} & \textbf{Agentic RAG} & \textbf{Agentic RAG} \\
\hline
\end{tabular}
\end{table}

A1 คือ RAG มาตรฐาน ไม่มีชั้นตัดสินใจและค้นเพียงรอบเดียว A3C มีชั้นเลือกการกระทำแต่ให้กฎในโปรแกรมเป็นผู้ตัดสินการค้นซ้ำ ส่วน A3 ให้แบบจำลองภาษาตัดสินทั้งสองชั้น

การแบ่งเช่นนี้ทำให้ได้คู่เปรียบเทียบที่ตอบคนละคำถาม คู่ A1 เทียบ A3C ตอบว่าการเพิ่มความเป็นเอเจนต์เข้าไปช่วยหรือไม่ แต่ต่างกัน 2 ตัวแปรพร้อมกัน ส่วนคู่ A3C เทียบ A3 เป็นคู่เดียวที่\textbf{เปลี่ยนตัวแปรทีละตัวอย่างแท้จริง} คือตอบว่าควรขยายขอบเขตการตัดสินใจของเอเจนต์ไปถึงการคุมการค้นซ้ำด้วยหรือไม่

\subsection{ลำดับการทำงานของระบบ}

ลำดับการทำงานเต็มรูปของ A3 มี 9 ขั้น โดย A3C ตัดขั้นที่ 7 ออกแล้วใช้กฎในโปรแกรมแทน ส่วน A1 ตัดทั้งขั้นที่ 2 และขั้นที่ 7 ถึง 8 ออก ดังนี้

\begin{enumerate}[label=\arabic*)]
    \item ปรับปรุงโปรไฟล์ความต้องการของผู้ใช้จากบทสนทนาที่ผ่านมา
    \item เลือกการกระทำของรอบนี้ ระหว่างถามกลับ สนทนา หรือค้นหา
    \item ประกอบคำค้นจากโปรไฟล์ความต้องการ
    \item แปลคำค้นเป็นภาษาอังกฤษหากผู้ใช้ถามเป็นภาษาไทย
    \item ค้นในคลังสินค้าโดยดึงผู้เข้าชิง 100 รายการ
    \item จัดอันดับผู้เข้าชิงและคัดด้วยเกณฑ์สัมบูรณ์
    \item ตัดสินว่าจะค้นซ้ำอีกรอบหรือไม่
    \item เขียนคำค้นใหม่สำหรับรอบถัดไป
    \item เขียนคำตอบให้ผู้ใช้จากสินค้าที่เลือกได้
\end{enumerate}

การเปรียบเทียบลำดับการทำงานของทั้งสามสถาปัตยกรรมแบบวางเคียงกันแสดงดังภาพที่ \ref{fig:arch} ซึ่งทำให้เห็นว่าชั้นที่ต่างกันมีเพียงขั้นที่ 2 และขั้นที่ 7 เท่านั้น

% =============================================================
% ภาพเปรียบเทียบลำดับการทำงานของทั้งสามสถาปัตยกรรม
% วาดด้วย TikZ ไม่ต้องใช้โปรแกรมภายนอก
% =============================================================

\begin{figure}[H]
\centering
\begin{tikzpicture}[
  every node/.style={transform shape},
  bx/.style   = {draw=black!55, rounded corners=1.5pt, minimum width=3.75cm,
                 minimum height=0.56cm, align=center, inner sep=1pt, font=\scriptsize},
  llm/.style  = {bx, fill=black!68, text=white},
  code/.style = {bx, fill=black!6},
  gone/.style = {bx, dashed, draw=black!35, text=black!45, fill=white},
  dec/.style  = {bx, fill=black!20, very thick, draw=black},
  hd/.style   = {minimum width=3.75cm, align=center, font=\small\bfseries},
  sn/.style   = {anchor=east, font=\scriptsize\itshape, text=black!55}
]

\def\cx{4.0}      % ระยะห่างระหว่างคอลัมน์
\def\ry{0.72}     % ระยะห่างระหว่างแถว

% ---------- หัวคอลัมน์ ----------
\node[hd] at (0,0)        {A1};
\node[hd] at (\cx,0)      {A3C};
\node[hd] at (2*\cx,0)    {A3};
\node[font=\scriptsize, text=black!55] at (0,-0.38)     {RAG มาตรฐาน};
\node[font=\scriptsize, text=black!55] at (\cx,-0.38)   {Agentic RAG};
\node[font=\scriptsize, text=black!55] at (2*\cx,-0.38) {Agentic RAG};

% ---------- เลขขั้น ----------
\foreach \i/\y in {1/1,2/2,3/3,4/4,5/5,6/6,7/7,8/8,9/9}
  \node[sn] at (-2.05,{-0.75-\y*\ry}) {\i};

% ---------- A1 ----------
\node[llm]  at (0,{-0.75-1*\ry}) {ปรับปรุงโปรไฟล์ผู้ใช้};
\node[gone] at (0,{-0.75-2*\ry}) {ไม่มี — ค้นทันทีทุกครั้ง};
\node[code] at (0,{-0.75-3*\ry}) {ประกอบคำค้นจากโปรไฟล์};
\node[llm]  at (0,{-0.75-4*\ry}) {แปลคำค้นเป็นอังกฤษ};
\node[code] at (0,{-0.75-5*\ry}) {ค้นในคลัง};
\node[code] at (0,{-0.75-6*\ry}) {จัดอันดับ + เกณฑ์คัด 0.55};
\node[gone] at (0,{-0.75-7*\ry}) {ไม่มี — จบรอบเดียว};
\node[gone] at (0,{-0.75-8*\ry}) {ไม่มี};
\node[llm]  at (0,{-0.75-9*\ry}) {เขียนคำตอบ};

% ---------- A3C ----------
\node[llm]  at (\cx,{-0.75-1*\ry}) {ปรับปรุงโปรไฟล์ผู้ใช้};
\node[dec]  at (\cx,{-0.75-2*\ry}) {LLM เลือก ถาม / คุย / ค้น};
\node[code] at (\cx,{-0.75-3*\ry}) {ประกอบคำค้นจากโปรไฟล์};
\node[llm]  at (\cx,{-0.75-4*\ry}) {แปลคำค้นเป็นอังกฤษ};
\node[code] at (\cx,{-0.75-5*\ry}) {ค้นในคลัง};
\node[code] at (\cx,{-0.75-6*\ry}) {จัดอันดับ + เกณฑ์คัด 0.55};
\node[dec]  at (\cx,{-0.75-7*\ry}) {กฎในโปรแกรม: ได้ครบ 3 ไหม};
\node[llm]  at (\cx,{-0.75-8*\ry}) {LLM เขียนคำค้นใหม่};
\node[llm]  at (\cx,{-0.75-9*\ry}) {เขียนคำตอบ};

% ---------- A3 ----------
\node[llm]  at (2*\cx,{-0.75-1*\ry}) {ปรับปรุงโปรไฟล์ผู้ใช้};
\node[dec]  at (2*\cx,{-0.75-2*\ry}) {LLM เลือก ถาม / คุย / ค้น};
\node[code] at (2*\cx,{-0.75-3*\ry}) {ประกอบคำค้นจากโปรไฟล์};
\node[llm]  at (2*\cx,{-0.75-4*\ry}) {แปลคำค้นเป็นอังกฤษ};
\node[code] at (2*\cx,{-0.75-5*\ry}) {ค้นในคลัง};
\node[code] at (2*\cx,{-0.75-6*\ry}) {จัดอันดับ + เกณฑ์คัด 0.55};
\node[dec]  at (2*\cx,{-0.75-7*\ry}) {\textbf{LLM ตัดสินเองว่าจะค้นซ้ำ}};
\node[llm]  at (2*\cx,{-0.75-8*\ry}) {LLM เขียนคำค้นใหม่};
\node[llm]  at (2*\cx,{-0.75-9*\ry}) {เขียนคำตอบ};

% ---------- ลูกศรวนกลับ ----------
\draw[-{Stealth[length=4pt]}, thick, black!70]
  ($(\cx+1.95,{-0.75-8*\ry})$) -- ++(0.42,0)
  |- ($(\cx+1.95,{-0.75-5*\ry})$);
\draw[-{Stealth[length=4pt]}, thick, black!70]
  ($(2*\cx+1.95,{-0.75-8*\ry})$) -- ++(0.42,0)
  |- ($(2*\cx+1.95,{-0.75-5*\ry})$);
\node[font=\tiny, text=black!60, rotate=90] at (\cx+2.62,{-0.75-6.5*\ry})   {ค้นซ้ำ};
\node[font=\tiny, text=black!60, rotate=90] at (2*\cx+2.62,{-0.75-6.5*\ry}) {ค้นซ้ำ};

% ---------- คำอธิบายสัญลักษณ์ ----------
\node[llm,  minimum width=1.5cm] (lg1) at (0.5,{-0.75-10.3*\ry})     {เรียก LLM};
\node[code, minimum width=1.5cm] (lg2) at (2.3,{-0.75-10.3*\ry})     {โปรแกรม};
\node[dec,  minimum width=2.3cm] (lg3) at (4.5,{-0.75-10.3*\ry})     {ชั้นที่ตัดสินใจ};
\node[gone, minimum width=1.6cm] (lg4) at (6.6,{-0.75-10.3*\ry})     {ไม่มีขั้นนี้};

\end{tikzpicture}
\caption{เปรียบเทียบลำดับการทำงาน 9 ขั้นของทั้งสามสถาปัตยกรรม}
\label{fig:arch}
\end{figure}


ขั้นที่ 2 มีผลสำคัญคือ หากเลือกถามกลับหรือสนทนา ระบบจะตอบทันทีโดย\textbf{ไม่ค้นเลย} และกำหนดเพดานการค้นซ้ำไว้ที่ 2 รอบทั้ง A3C และ A3

% -----------------------------------------------------------
\section{การออกแบบส่วนประกอบของระบบ}
% -----------------------------------------------------------

ส่วนประกอบทั้งหมดใช้ร่วมกันทุกสถาปัตยกรรม ดังตารางที่ \ref{tab:modules}

\begin{table}[H]
\centering
\caption{ส่วนประกอบของระบบและค่าที่ตั้งไว้}
\label{tab:modules}
\small
\begin{tabular}{|l|p{4.1cm}|p{3.3cm}|c|c|}
\hline
\textbf{ส่วนประกอบ} & \textbf{หน้าที่} & \textbf{ค่าที่ตั้งไว้} & \textbf{A1} & \textbf{A3C} \\
\hline
ความจำ & จดจำความต้องการข้ามรอบสนทนาเป็นโปรไฟล์ 9 ช่องที่สะสมขึ้นเรื่อย ๆ & อ่านย้อนหลัง 4 รอบ ขยายงบได้ไม่เกิน 2 ครั้ง & \cmark & \cmark \\ \hline
วางแผน & เลือกการกระทำของรอบนี้ & ถามกลับได้ไม่เกิน 2 ครั้ง มีทางสำรองเป็นกฎเสมอ & -- & \cmark \\ \hline
แปลภาษา & แปลคำค้นไทยเป็นอังกฤษก่อนค้น & ค่าความสุ่มเป็นศูนย์ จำคำแปลไว้ไม่เรียกซ้ำ & \cmark & \cmark \\ \hline
ค้นคืน & ค้นผู้เข้าชิงจากคลัง & ดึง 100 รายการต่อการค้นหนึ่งครั้ง & \cmark & \cmark \\ \hline
จัดอันดับ & จัดอันดับและคัดด้วยเกณฑ์สัมบูรณ์ & ส่งต่อ 3 รายการ เกณฑ์คัด 0.55 & \cmark & \cmark \\ \hline
เขียนคำค้นใหม่ & สร้างคำค้นสำหรับรอบถัดไป & ใช้แบบจำลองภาษา ไม่ใช้กฎตัดคำ & -- & \cmark \\ \hline
สร้างคำตอบ & เขียนคำตอบให้ผู้ใช้ & ค่าความสุ่มเป็นศูนย์ ส่งเฉพาะข้อเท็จจริงจากคลัง & \cmark & \cmark \\
\hline
\end{tabular}
\end{table}

ส่วนเขียนคำค้นใหม่เลือกใช้แบบจำลองภาษาแทนการเขียนกฎ เนื่องจากรุ่นแรกของระบบใช้กฎตัดคำขยายออกทีละชั้น แล้วพบว่า\textbf{กฎตัดชื่อสินค้าทิ้ง} ทำให้คำค้นรอบถัดไปกว้างกว่าเดิมและได้สินค้าที่ผิดมากขึ้น

% -----------------------------------------------------------
\section{การออกแบบชั้นจัดอันดับ}
% -----------------------------------------------------------

% ลำดับขั้นย่อยภายในชั้นจัดอันดับ — แทรกก่อนหัวข้อการรวมสัญญาณในบทที่ 4

\subsection{ลำดับการทำงานภายในชั้นจัดอันดับ}

ชั้นจัดอันดับรับผู้เข้าชิง 100 รายการจากขั้นค้นคืน แล้วทำงาน 4 ขั้นย่อยตามลำดับ
ดังตารางที่ \ref{tab:rankorder}

\begin{table}[H]
\centering
\caption{ลำดับขั้นย่อยภายในชั้นจัดอันดับ พร้อมจำนวนที่เหลือจากตัวอย่างการรันจริง}
\label{tab:rankorder}
\small
\begin{tabular}{|c|p{7.2cm}|c|}
\hline
\textbf{ขั้นย่อย} & \textbf{การทำงาน} & \textbf{จำนวนที่เหลือ} \\
\hline
1 & ตัดทิ้งตามเงื่อนไขที่ผู้ใช้ระบุเอง เช่น ราคาเกินงบ หรือแบรนด์ที่ผู้ใช้ปฏิเสธ & 100 \\ \hline
2 & คัดด้วยเกณฑ์สัมบูรณ์ 0.55 โดยใช้คะแนนความใกล้เชิงความหมายเพียงอย่างเดียว & 73 \\ \hline
3 & \textbf{จัดอันดับด้วย 5 สัญญาณ รวมถึงการนับคำซ้อน} & 73 \\ \hline
4 & หยิบ 3 รายการแรก โดยเลี่ยงการหยิบจากร้านค้าเดียวกันซ้ำ & 3 \\
\hline
\end{tabular}
\end{table}

\textbf{ประเด็นที่ต้องเข้าใจให้ตรงคือการนับคำซ้อนไม่ได้ทำหน้าที่คัดสินค้าทิ้ง}
สัญญาณนี้ทำงานที่ขั้นย่อยที่ 3 ซึ่งอยู่หลังการคัดด้วยเกณฑ์สัมบูรณ์แล้ว
หน้าที่ของมันคือเรียงลำดับสินค้าที่ผ่านเกณฑ์มาแล้วเท่านั้น
สินค้าที่ได้คะแนนการนับคำซ้อนเป็นศูนย์จึงยังอยู่ในกอง เพียงแต่ถูกจัดไปอยู่อันดับท้าย

การแยกหน้าที่เช่นนี้จำเป็น เพราะการคัดว่าดีพอหรือยังต้องใช้เกณฑ์สัมบูรณ์
ส่วนการจัดอันดับเป็นการเปรียบเทียบกันเองภายในกอง ถ้านำการนับคำซ้อนไปใช้เป็นเกณฑ์คัด
สินค้าที่คลังเรียกด้วยคำอื่นจะถูกตัดทิ้งทั้งที่ตรงกับความต้องการ
ซึ่งเป็นข้อผิดพลาดที่พบจริงและรายงานไว้ในหัวข้อผลเชิงลบ

ข้อความที่นำมานับคำซ้อนคือคำค้นหลังผ่านขั้นแปลภาษาแล้ว ไม่ใช่ข้อความที่ผู้ใช้พิมพ์โดยตรง
ฝั่งภาษาไทยจึงได้รับผลกระทบจากความคลาดเคลื่อนของขั้นแปลอย่างหลีกเลี่ยงไม่ได้

นอกจากนี้ ชั้นจัดอันดับถูกเรียกใช้ทุกครั้งที่มีการค้น ไม่ใช่เพียงครั้งเดียวต่อหนึ่งคำถาม
เมื่อการค้นซ้ำทำงาน ผู้เข้าชิงชุดใหม่จะถูกจัดอันดับใหม่ทั้งหมด
และชั้นนี้เปิดช่องให้เปลี่ยนน้ำหนักของการนับคำซ้อนขณะทำงานได้
ซึ่งเป็นช่องทางที่ใช้ทดลองให้เอเจนต์ปรับน้ำหนักเอง และพบว่าไม่ได้ผล


\subsection{การรวมสัญญาณด้วยการรวมอันดับ}

ชั้นจัดอันดับรวมสัญญาณหลายตัวด้วยวิธีรวมอันดับตามสมการที่ \ref{eq:rrf} โดยใช้ค่าคงที่ $k = 60$ สัญญาณและน้ำหนักที่ใช้แสดงดังตารางที่ \ref{tab:signals}

\begin{table}[H]
\centering
\caption{สัญญาณที่ใช้จัดอันดับและน้ำหนัก}
\label{tab:signals}
\begin{tabular}{|c|l|c|p{6.4cm}|}
\hline
\textbf{ลำดับ} & \textbf{สัญญาณ} & \textbf{น้ำหนัก} & \textbf{เหตุผล} \\
\hline
1 & ความใกล้เชิงความหมาย & 3 & สัญญาณหลักที่บอกความตรงของเนื้อหา \\ \hline
2 & คะแนนดาวเฉลี่ย & 1 & คุณภาพสินค้าในสายตาผู้ซื้อ \\ \hline
3 & จำนวนรีวิว & 1 & กันสินค้า 5 ดาวรีวิวเดียวชนะสินค้า 4.8 ดาวพันรีวิว \\ \hline
4 & ความตรงของหมวดสินค้า & 1 & ตรงกับหมวดที่ผู้ใช้ระบุในโปรไฟล์หรือไม่ \\ \hline
5 & \textbf{จำนวนคำที่ซ้อนกับคำค้น} & \textbf{6} & \textbf{สัญญาณที่เพิ่มเข้ามาภายหลัง} จับคำเงื่อนไขเฉพาะที่สัญญาณอื่นจับไม่ได้ \\
\hline
\end{tabular}
\end{table}

\subsection{เกณฑ์คัดแบบสัมบูรณ์}

ชั้นจัดอันดับคืนค่า 3 อย่าง คือรายการที่เลือกแล้ว รายการทั้งหมดที่ผ่านเกณฑ์พร้อมเหตุผล และ\textbf{จำนวนรายการที่ผ่านเกณฑ์} ค่าสุดท้ายนี้เป็นสิ่งที่การค้นซ้ำใช้ตัดสินใจ

เหตุผลที่ต้องมีเกณฑ์สัมบูรณ์แยกจากการจัดอันดับคือ\textbf{อันดับเป็นการเปรียบเทียบของในกองเดียวกัน แต่คำถามว่าดีพอหรือยังต้องใช้เกณฑ์สัมบูรณ์} สินค้าที่แย่ที่สุดใน 100 รายการก็ยังได้อันดับที่ 100 อยู่ดี การจัดอันดับจึงไม่เคยบอกได้ว่าทั้งกองไม่ตรงเลย

ค่าเกณฑ์คัดกำหนดไว้ที่ 0.55 ซึ่งได้จากการสร้างกองคำค้นแยกต่างหากที่รู้คำตอบล่วงหน้าว่าควรผ่านหรือควรตก แล้วหาค่าที่แยกสองกองได้ดีที่สุด กองนี้แยกจากชุดทดสอบที่ใช้รายงานผลเพื่อไม่ให้เกิดการปรับค่าบนข้อมูลที่ใช้วัดผล

% -----------------------------------------------------------
\section{การสร้างชุดทดสอบและการติดเฉลย}
% -----------------------------------------------------------

\subsection{ที่มาของคำค้น}

ชุดทดสอบสร้างจากคำค้นจริงของผู้ใช้ที่ยืมมาจากชุดข้อมูลมาตรฐาน 2 ชุด \cite{reddy2022esci, chen2022wands} นำมายิงเข้าคลังของโครงงานนี้จำนวน 1{,}980 คำ แล้วคัดด้วยเกณฑ์ตัวเลขเหลือ 140 คำค้น ดังตารางที่ \ref{tab:testset}

\begin{table}[H]
\centering
\caption{องค์ประกอบของชุดทดสอบ 140 คำค้น}
\label{tab:testset}
\small
\begin{tabular}{|l|c|p{4.3cm}|p{3.1cm}|}
\hline
\textbf{กลุ่ม} & \textbf{จำนวน} & \textbf{เกณฑ์ที่ใช้คัด} & \textbf{คำตอบที่ถูก} \\
\hline
คลังมีสินค้าตรง & 50 & คะแนนอันดับ 1 ไม่ต่ำกว่า 0.50 และผ่านเกณฑ์คัดอย่างน้อย 3 รายการ & แนะนำสินค้า 3 รายการที่ตรง \\ \hline
คลังไม่มีสินค้าตรง & 60 & คะแนนอันดับ 1 ต่ำกว่า 0.47 & ตอบว่าไม่พบ \\ \hline
มีเงื่อนไขงบประมาณ & 30 & นำคำค้นกลุ่มแรกมาเติมเพดานราคา & แนะนำสินค้าที่ไม่เกินงบ \\
\hline
\end{tabular}
\end{table}

เหตุผลที่ไม่เขียนคำค้นขึ้นเองคือได้ทดลองเขียนเองแล้ว 2 รอบ รวม 76 คำ แต่\textbf{ได้คำค้นชั้นยาก 0 คำทั้งสองรอบ} เพราะผู้เขียนไม่รู้รหัสอะไหล่จริง ศัพท์เฉพาะทาง และคำที่ผู้ใช้จริงสะกดผิด ส่วนเหตุผลที่ไม่ยืมเฉลยของชุดข้อมูลต้นทางมาด้วยคือคลังสินค้าเป็นคนละคลังกัน โดยคลังต้นทางซ้อนทับคลังของโครงงานนี้เพียง 17 จาก 856{,}238 รายการเมื่อเทียบด้วยรหัสสินค้า เฉลยของต้นทางจึงใช้ไม่ได้

\subsection{การติดเฉลยและการป้องกันอคติ}

เฉลยความตรงติดตามเกณฑ์ 4 ระดับ \cite{reddy2022esci} คือตรงเป๊ะ ใช้แทนกันได้ ของเสริม และไม่เกี่ยว โดยผู้วิจัยติดเองทั้งหมด เฉลยที่ใช้คิดตัวเลขในรายงานนี้มีจำนวน 424 คู่

มาตรการป้องกันอคติที่ใช้มี 3 ข้อ

\begin{enumerate}[label=\arabic*)]
    \item \textbf{ปิดตาผู้ติดป้าย} ไฟล์เฉลยไม่ระบุว่าสถาปัตยกรรมใดเป็นผู้หยิบสินค้ารายการนั้นขึ้นมา ไม่มีคอลัมน์รอบ อันดับ หรือคะแนน และเรียงลำดับแถวด้วยค่าแฮชคงที่ หากผู้ติดป้ายรู้ว่าสินค้ามาจากสถาปัตยกรรมใด เฉลยย่อมเอนไปตามความคาดหวังว่าแบบใดควรชนะ ซึ่งเป็นสิ่งที่โครงงานนี้กำลังวัดอยู่พอดี
    \item \textbf{ล็อกชุดคำค้นก่อนรัน} และห้ามแก้หลังเห็นผล
    \item \textbf{แยกกองพัฒนาออกจากกองวัดผล} คำค้นที่ใช้เขียนและปรับคำสั่งของแบบจำลองห้ามนำมาใช้รายงานผล
\end{enumerate}

\subsection{ความเอนเอียงของกองที่ติดเฉลย}

เมื่อเปลี่ยนชั้นจัดอันดับ พบว่าคำค้นเดิมและกองผู้เข้าชิงเดิม แต่\textbf{สามอันดับแรกเปลี่ยนไปถึง 51 จาก 80 คำค้น} สินค้าที่โผล่ขึ้นมาใหม่ยังไม่มีเฉลย และวิธีคิดตัวชี้วัดตัดคำค้นที่เฉลยไม่ครบทิ้งทั้งข้อ ผลคือชุดคำค้นที่เทียบได้ครบทั้งสามสถาปัตยกรรมเหลือเพียง 21 จาก 50 คำ และใน 21 คำนั้นมี 19 คำที่ทุกสถาปัตยกรรมให้ผลเหมือนกันทุกประการ กล่าวคือ\textbf{เหลือเฉพาะคำค้นที่ไม่มีอะไรเปลี่ยน จึงวัดอะไรไม่ได้เลย}

ปรากฏการณ์นี้คือความเอนเอียงของกองที่ติดเฉลย \cite{zobel1998reliable} ทางแก้คือติดเฉลยเพิ่มสำหรับสินค้าที่โผล่ใหม่ ซึ่งได้ดำเนินการแล้วจนชุดคำค้นที่เทียบได้ครบกลับมาเป็น 140 จาก 140 คำ

% -----------------------------------------------------------
\section{ผลการทดลอง}
% -----------------------------------------------------------

\subsection{ผลรวมทุกตัวชี้วัด}

ผลการวัดทั้งหมดได้จากการรัน 3 สถาปัตยกรรม 2 ภาษา ภาษาละ 1 รอบ รวม 840 คำตอบ ไม่มีข้อใดทำงานผิดพลาด แสดงดังตารางที่ \ref{tab:results_all}

\begin{table}[H]
\centering
\caption{ผลการวัดทุกตัวชี้วัด แยกตามสถาปัตยกรรมและภาษา}
\label{tab:results_all}
\small
\begin{tabular}{|p{4.4cm}|c|c|c|c|c|c|}
\hline
\textbf{ตัวชี้วัด} & \textbf{A1 อ.} & \textbf{A3C อ.} & \textbf{A3 อ.} & \textbf{A1 ท.} & \textbf{A3C ท.} & \textbf{A3 ท.} \\
\hline
nDCG@3 ภาพรวม 140 คำ & 0.3288 & \textbf{0.3564} & \textbf{0.3564} & 0.3165 & \textbf{0.3267} & \textbf{0.3267} \\ \hline
nDCG@3 คลังมีของ & 0.594 & 0.594 & 0.594 & 0.545 & \textbf{0.549} & \textbf{0.549} \\ \hline
nDCG@3 มีเงื่อนไขราคา & 0.544 & \textbf{0.549} & \textbf{0.549} & 0.473 & \textbf{0.492} & \textbf{0.492} \\ \hline
nDCG@3 คลังไม่มีของ & 0.000 & \textbf{0.062} & \textbf{0.062} & 0.048 & \textbf{0.059} & \textbf{0.059} \\ \hline
P@3 ตรงหรือใช้แทนได้ & 64.0\% & 64.0\% & 64.0\% & 55.3\% & 55.3\% & 55.3\% \\ \hline
ได้สินค้าครบ 3 รายการ & 79 & \textbf{86} & \textbf{86} & 78 & \textbf{83} & \textbf{83} \\ \hline
ปฏิเสธถูกเมื่อคลังไม่มีของ & \textbf{60/60} & 50/60 & 50/60 & \textbf{55/60} & 51/60 & 51/60 \\ \hline
ตัดสินถูกเมื่อนับสองทิศ & 54/60 & \textbf{56/60} & \textbf{56/60} & 58/60 & 58/60 & 58/60 \\ \hline
จำนวนครั้งที่เรียกแบบจำลอง & \textbf{1.56} & 3.36 & 4.68 & \textbf{2.06} & 3.86 & 5.19 \\ \hline
เวลาต่อข้อ (วินาที) & \textbf{7.5} & 14.2 & 19.6 & \textbf{9.6} & 16.0 & 21.3 \\
\hline
\end{tabular}
\end{table}

\textbf{หมายเหตุสำคัญเรื่องการอ่านตารางนี้} แถวปฏิเสธถูกนับการปฏิเสธทุกครั้งว่าถูก จึงให้คะแนนเต็มกับสถาปัตยกรรมที่ปฏิเสธทุกอย่าง ส่วนแถวตัดสินถูกนับทั้งสองทิศ คือปฏิเสธถูกเมื่อคลังไม่มีของ บวกกับแนะนำได้เมื่อคลังมีของ เมื่อนับแบบนี้ A1 ฝั่งอังกฤษตกจาก 60/60 เหลือ 54/60 ซึ่งแสดงว่าคะแนนเต็มเดิมมาจากการปฏิเสธทุกอย่าง ไม่ใช่การตัดสินที่ถูกต้อง

\subsection{การทดสอบนัยสำคัญทางสถิติ}

ผลการทดสอบจับคู่แบบสลับเครื่องหมายที่นับครบทุกกรณีแสดงดังตารางที่ \ref{tab:signif}

\begin{table}[H]
\centering
\caption{ผลการทดสอบนัยสำคัญของ A3C และ A3 เทียบกับ A1}
\label{tab:signif}
\begin{tabular}{|l|c|c|c|c|}
\hline
\textbf{ภาษา} & \textbf{A1} & \textbf{A3C และ A3} & \textbf{ดีขึ้น $\cdot$ แย่ลง} & \textbf{ค่า $p$} \\
\hline
อังกฤษ & 0.3288 & \textbf{0.3564} & \textbf{8 $\cdot$ 0} & \textbf{0.0078} \\ \hline
ไทย & 0.3165 & \textbf{0.3267} & \textbf{5 $\cdot$ 0} & 0.0625 \\
\hline
\end{tabular}
\end{table}

ฝั่งอังกฤษมีนัยสำคัญที่ระดับ 0.01 ส่วนฝั่งไทยได้ค่า 0.0625 ซึ่ง\textbf{เป็นค่าต่ำสุดที่เป็นไปได้เมื่อมีคำค้นที่ให้ผลต่างกันเพียง 5 ข้อ} เนื่องจาก $2/2^5 = 0.0625$ จึงเป็นข้อจำกัดของจำนวนข้อ ไม่ใช่สัญญาณว่าผลอ่อน ประเด็นสำคัญคือ\textbf{ทั้งสองภาษาไม่มีคำค้นใดที่สถาปัตยกรรมแบบเอเจนต์แพ้เลยแม้แต่ข้อเดียว}

\subsection{ผลแยกตามกลุ่มคำค้น}

ผลแยกตามกลุ่มแสดงดังตารางที่ \ref{tab:bygroup} ซึ่งชี้ว่าความได้เปรียบของสถาปัตยกรรมแบบเอเจนต์กระจุกอยู่ที่กลุ่มเดียว

\begin{table}[H]
\centering
\caption{ผลแยกตามกลุ่มคำค้นและภาษา}
\label{tab:bygroup}
\begin{tabular}{|l|l|c|c|c|c|}
\hline
\textbf{กลุ่ม} & \textbf{ภาษา} & \textbf{A1} & \textbf{A3C และ A3} & \textbf{ดีขึ้น $\cdot$ แย่ลง} & \textbf{ค่า $p$} \\
\hline
คลังมีของ & อังกฤษ & 0.5944 & 0.5944 & 0 $\cdot$ 0 & 1.0000 \\ \hline
คลังมีของ & ไทย & 0.5454 & 0.5490 & 1 $\cdot$ 0 & 1.0000 \\ \hline
มีเงื่อนไขราคา & อังกฤษ & 0.5438 & 0.5490 & 1 $\cdot$ 0 & 1.0000 \\ \hline
มีเงื่อนไขราคา & ไทย & 0.4726 & 0.4922 & 2 $\cdot$ 0 & 0.5000 \\ \hline
\textbf{คลังไม่มีของ} & \textbf{อังกฤษ} & \textbf{0.0000} & \textbf{0.0617} & \textbf{7 $\cdot$ 0} & \textbf{0.0156} \\ \hline
คลังไม่มีของ & ไทย & 0.0477 & 0.0588 & 2 $\cdot$ 0 & 0.5000 \\
\hline
\end{tabular}
\end{table}

\textbf{ข้อสังเกตที่สำคัญที่สุดจากตารางนี้} คือในกลุ่มที่คลังมีสินค้าตรง ทั้งสามสถาปัตยกรรมได้คะแนนเท่ากันทุกทศนิยมในฝั่งอังกฤษ แปลว่าเมื่อการค้นรอบแรกได้ผลดีอยู่แล้ว ชั้นตัดสินใจไม่ได้สร้างความต่างใด ๆ ความได้เปรียบทั้งหมดมาจากกลุ่มที่คลังไม่มีสินค้าตรง ซึ่งการค้นซ้ำทำให้ระบบหาสินค้าที่ใช้แทนกันได้เจอ ในคำค้นที่ A1 ตอบว่าไม่พบ

สถิติการทำงานของการค้นซ้ำสนับสนุนข้อสังเกตนี้ คือการค้นซ้ำ\textbf{ไม่ทำงานเลยเมื่อคลังมีของ} และทำงาน 36 เปอร์เซ็นต์ฝั่งอังกฤษกับ 38 เปอร์เซ็นต์ฝั่งไทย ซึ่งอยู่ในกลุ่มที่คลังไม่มีของทั้งหมด

% -----------------------------------------------------------
\section{ข้อค้นพบ}
% -----------------------------------------------------------

\subsection{ชั้นจัดอันดับเป็นคอขวดของระบบ ไม่ใช่ชั้นค้นคืน}

เพื่อหาว่าสินค้าที่ดีสูญหายไปที่ขั้นใด ได้ตรวจสอบทุกขั้นของระบบในกลุ่มที่คลังมีสินค้า โดยนับสินค้าที่มีเฉลยระดับใช้แทนได้ขึ้นไปว่าตกหล่นที่ชั้นใด ผลแสดงดังตารางที่ \ref{tab:funnel}

\begin{table}[H]
\centering
\caption{ชั้นที่ทำให้สินค้าที่ดีสูญหาย}
\label{tab:funnel}
\begin{tabular}{|l|c|c|}
\hline
\textbf{ชั้นที่สูญหาย} & \textbf{ไทย} & \textbf{อังกฤษ} \\
\hline
ค้นไม่เจอใน 100 อันดับแรก & 16 (5.9\%) & 18 (6.6\%) \\ \hline
ตกที่กฎกรองทิ้ง & 0 & 0 \\ \hline
ตกที่เกณฑ์คัด & 16 (5.9\%) & 17 (6.2\%) \\ \hline
\textbf{ตกที่ชั้นจัดอันดับ} & \textbf{131 (48.2\%)} & \textbf{123 (45.2\%)} \\ \hline
ตกที่การตัดร้านค้าซ้ำ & 3 & 2 \\
\hline
\end{tabular}
\end{table}

\textbf{ชั้นจัดอันดับทำสินค้าที่ดีสูญหายมากกว่าทุกชั้นอื่นรวมกันประมาณ 3.5 เท่า และสูญหายในสัดส่วนใกล้เคียงกันทั้งสองภาษา} ข้อค้นพบนี้เปลี่ยนทิศทางการพัฒนา เพราะก่อนหน้านี้ความพยายามปรับปรุงมุ่งไปที่ชั้นค้นคืนและชั้นตัดสินใจ ซึ่งไม่ใช่จุดที่เสียหายจริง

\subsection{สัญญาณที่ขาดไปของชั้นจัดอันดับ}

เมื่อตรวจว่าสัญญาณเดิม 4 ตัวแยกสินค้าที่ตรงเป๊ะออกจากสินค้าที่ไม่เกี่ยวได้หรือไม่ พบว่าแยกไม่ได้ ดังตารางที่ \ref{tab:signal_power}

\begin{table}[H]
\centering
\caption{ความสามารถของแต่ละสัญญาณในการแยกสินค้าที่ตรงออกจากสินค้าที่ไม่ตรง}
\label{tab:signal_power}
\begin{tabular}{|l|p{7.6cm}|}
\hline
\textbf{สัญญาณ} & \textbf{ผลการแยก} \\
\hline
ความใกล้เชิงความหมาย & แยกไม่ได้ คะแนนเฉลี่ย 0.639 เทียบ 0.591 แทบไม่ต่างกัน \\ \hline
ดาว จำนวนรีวิว และหมวดสินค้า & แยกไม่ได้ เพราะไม่ได้วัดความตรงเลย \\ \hline
\textbf{จำนวนคำที่ซ้อนกับคำค้น} & \textbf{แยกได้ 0.667 เทียบ 0.333} \\
\hline
\end{tabular}
\end{table}

สาเหตุที่คะแนนความใกล้เชิงความหมายแยกไม่ได้ เกิดจากเวกเตอร์ของข้อความทั้งหมดกระจุกอยู่ในกรวยแคบ ซึ่งเป็นคุณสมบัติที่พบในแบบจำลองภาษาโดยทั่วไป \cite{ethayarajh2019contextual} การวัดพบว่าข้อความที่ไม่เกี่ยวข้องกันเลยยังได้คะแนนประมาณ 0.457 ไม่ใช่ศูนย์ ช่วงค่าที่ใช้งานได้จริงจึงแคบเพียง 0.42 ถึง 0.75

เมื่อเพิ่มการนับคำที่ซ้อนกันเป็นสัญญาณที่ 5 ได้ผลเพิ่มขึ้น \textbf{4.5 จุด} บนกองที่จำกัดเฉพาะสินค้าที่มีเฉลย ที่ระดับนัยสำคัญ $p = 0.0005$ และยืนยันผลแบบตลอดสายที่ A1 ฝั่งอังกฤษ ได้เพิ่มขึ้น 3.5 จุดในกลุ่มคลังมีของ และ 3.4 จุดในกลุ่มมีเงื่อนไขราคา ซึ่งเป็นผลที่\textbf{มากกว่าผลที่ความเป็นเอเจนต์ให้ประมาณ 2 เท่า}

% วิธีนับคำซ้อนและกฎตัดคำ — แทรกในบทที่ 4 หัวข้อสัญญาณที่ขาดไปของชั้นจัดอันดับ

\subsubsection*{วิธีนับคำซ้อนและกฎตัดคำ}

การนับทำเป็น 5 ขั้น คือตัดข้อความทั้งของคำค้นและของชื่อสินค้าให้เป็นคำ
ทิ้งคำที่ไม่ช่วยแยก เอาเฉพาะคำที่ไม่ซ้ำ นับว่าคำในคำค้นปรากฏในชื่อสินค้ากี่คำ
แล้วหารด้วยจำนวนคำในคำค้น ได้ค่าระหว่าง 0 ถึง 1
ตัวหารเป็นจำนวนคำของคำค้น ไม่ใช่ของชื่อสินค้า สินค้าที่มีชื่อยาวจึงไม่ได้เปรียบ

ลำดับการทำงานทั้งหมดแสดงดังภาพที่ \ref{fig:wordcount}

% =============================================================
% ลำดับการทำงานของสัญญาณนับคำซ้อน
% =============================================================

\begin{figure}[H]
\centering
\begin{tikzpicture}[
  bx/.style  = {draw=black!55, rounded corners=2pt, minimum width=3.3cm,
                minimum height=0.8cm, align=center, inner sep=2pt, font=\scriptsize},
  inp/.style = {bx, fill=black!6},
  op/.style  = {bx},
  mrg/.style = {bx, fill=black!20, very thick, draw=black, minimum width=4.4cm},
  fin/.style = {bx, fill=black!68, text=white, minimum width=4.4cm},
  note/.style= {draw=black!40, dashed, rounded corners=2pt, align=left,
                inner sep=5pt, font=\scriptsize, text width=4.4cm, fill=white},
  ar/.style  = {-{Stealth[length=4pt]}, thick, black!65}
]

\def\lx{-2.3} \def\rx{2.3}

% ---------- ข้อมูลเข้า ----------
\node[inp] (q0) at (\lx, 0)      {คำค้น\\(หลังผ่านขั้นแปลภาษา)};
\node[inp] (t0) at (\rx, 0)      {ชื่อสินค้าที่ค้นได้};

% ---------- ตัดเป็นคำ ----------
\node[op]  (q1) at (\lx,-1.35)   {ตัดเป็นคำ\\เฉพาะ a--z และ 0--9};
\node[op]  (t1) at (\rx,-1.35)   {ตัดเป็นคำ\\เฉพาะ a--z และ 0--9};

% ---------- ทิ้งคำ ----------
\node[op]  (q2) at (\lx,-2.75)   {ทิ้งคำที่ไม่ช่วยแยก};
\node[op]  (t2) at (\rx,-2.75)   {ทิ้งคำที่ไม่ช่วยแยก};

% ---------- เซต ----------
\node[op]  (q3) at (\lx,-4.1)    {เซตคำของคำค้น\\$Q$};
\node[op]  (t3) at (\rx,-4.1)    {เซตคำของชื่อสินค้า\\$T$};

% ---------- รวม ----------
\node[mrg] (m)  at (0,-5.5)      {หาคำที่อยู่ทั้งสองฝั่ง\\$|Q \cap T|$};
\node[fin] (f)  at (0,-6.9)      {คะแนน $= \dfrac{|Q \cap T|}{|Q|}$ \; มีค่า 0 ถึง 1};

\foreach \i/\j in {q0/q1, q1/q2, q2/q3, t0/t1, t1/t2, t2/t3} \draw[ar] (\i) -- (\j);
\draw[ar] (q3) |- (m.west);
\draw[ar] (t3) |- (m.east);
\draw[ar] (m)  -- (f);

% ---------- กฎที่ใช้ทิ้งคำ ----------
\node[note] (n) at (7.2,-2.75)
  {\textbf{กฎทิ้งคำ 2 ข้อ}\\[2pt]
   1. อยู่ในรายการคำที่ไม่ช่วยแยก 13 คำ\\
   \hspace*{0.9em}คำไวยากรณ์ และคำที่พบบ่อย\\
   \hspace*{0.9em}ในชื่อสินค้าจนไม่บอกความตรง\\[3pt]
   2. สั้นกว่า 3 ตัวอักษร\\
   \hspace*{0.9em}เพราะตรงกันโดยบังเอิญได้ง่าย};
\draw[ar, dashed, black!45] (n.west) -- (t2.east);

% ---------- หมายเหตุตัวหาร ----------
\node[note, text width=4.4cm] (d) at (7.2,-6.4)
  {\textbf{ตัวหารคือ $|Q|$}\\[2pt]
   นับจากจำนวนคำของ\textbf{คำค้น}\\
   ไม่ใช่ของชื่อสินค้า\\
   สินค้าที่ชื่อยาวจึงไม่ได้เปรียบ};
\draw[ar, dashed, black!45] (d.west) -- (f.east);

\end{tikzpicture}
\caption{ลำดับการทำงานของสัญญาณนับคำซ้อน}
\label{fig:wordcount}
\end{figure}


ตัวอย่างการคำนวณแสดงดังตารางที่ \ref{tab:wordcount_example}

\begin{table}[H]
\centering
\caption{ตัวอย่างการนับคำซ้อนสำหรับคำค้น squeaky balls for large dogs}
\label{tab:wordcount_example}
\small
\begin{tabular}{|c|p{10.6cm}|}
\hline
\textbf{ขั้น} & \textbf{ผลลัพธ์} \\
\hline
1 & ตัดคำค้นเป็นคำ ได้ squeaky $\cdot$ balls $\cdot$ for $\cdot$ large $\cdot$ dogs \\ \hline
2 & ทิ้งคำที่ไม่ช่วยแยก เหลือ squeaky $\cdot$ balls $\cdot$ large $\cdot$ dogs รวม 4 คำ \\ \hline
3 & ชื่อสินค้า XL Thick Dog Rope Toy for Large Breed Aggressive Chewers \\ \hline
4 & คำที่ปรากฏในชื่อสินค้า มี large เพียงคำเดียว \\ \hline
5 & คะแนน $= 1 \div 4 = 0.250$ \\
\hline
\end{tabular}
\end{table}

แนวคิดของการวัดความตรงจากคำที่ตรงกันเป็นรากฐานของการค้นคืนแบบจับคู่คำ ซึ่งใช้กันมานานและยังเป็นเส้นฐานที่แข็งแรงจนถึงปัจจุบัน \cite{robertson2009bm25}\
กฎที่ใช้ตัดคำมี 2 ข้อ ข้อแรกคือตัดคำที่อยู่ในรายการคำที่ไม่ช่วยแยกจำนวน 13 คำ
ซึ่งแบ่งเป็นคำไวยากรณ์ที่ปรากฏในชื่อสินค้าแทบทุกชิ้น และคำที่พบบ่อยในชื่อสินค้า
จนไม่บอกความตรง เช่นคำที่หมายถึงชุด แพ็ก นิ้ว และขนาด
ข้อที่สองคือตัดคำที่สั้นกว่า 3 ตัวอักษร เพราะจับคู่ตรงกันโดยบังเอิญได้ง่าย

\textbf{ที่มาของกฎทั้งสองข้อคือการกำหนดด้วยมือ ไม่ได้มาจากการเรียนรู้จากข้อมูล}
และยังไม่ได้กวาดค่าเปรียบเทียบกับทางเลือกอื่น ซึ่งต่างจากน้ำหนักของสัญญาณ
และค่าเกณฑ์คัดที่เลือกมาจากการวัด ทางเลือกที่เป็นมาตรฐานกว่าคือการถ่วงน้ำหนักคำ
ตามความถี่ผกผันในคลัง ซึ่งจะทำให้ระบบหาเองได้ว่าคำใดไม่ช่วยแยกและปรับตามโดเมนได้เอง
โครงงานนี้เลือกรายการตายตัวเพราะอธิบายได้ทุกคำและให้ผลเดิมทุกครั้ง
ซึ่งจำเป็นกับงานที่ต้องทำซ้ำได้ แต่เป็นการแลกเปลี่ยน ไม่ใช่ทางเลือกที่ดีกว่าเสมอไป

\textbf{ผลข้างเคียงที่วัดได้ของกฎตัดคำสั้น} คือคำค้นที่เขียนหน่วยแยกจากตัวเลข
จะถูกตัดทั้งตัวเลขและหน่วยทิ้ง ทำให้สัญญาณนี้แยกความจุที่ต่างกันไม่ได้
ดังตัวอย่างในตารางที่ \ref{tab:wordcount_limit}

\begin{table}[H]
\centering
\caption{ผลของกฎตัดคำสั้นต่อคำค้นที่มีตัวเลขและหน่วย}
\label{tab:wordcount_limit}
\small
\begin{tabular}{|l|c|c|}
\hline
\textbf{รูปแบบคำค้น} & \textbf{สินค้า 2 TB} & \textbf{สินค้า 1 TB} \\
\hline
เขียนติดกัน \quad usb flash drive 2tb & 1.00 & \textbf{0.75} \\ \hline
เขียนแยกคำ \quad usb flash drive 2 tb & 1.00 & \textbf{1.00} \\
\hline
\end{tabular}
\end{table}

กล่าวคือสัญญาณนี้จับคำเงื่อนไขที่ยาวตั้งแต่ 3 ตัวอักษรขึ้นไปได้ดี
เช่น คำบอกลักษณะหรือชื่อแบรนด์ แต่จับตัวเลขและหน่วยสั้นที่เขียนแยกคำไม่ได้
ซึ่งเป็นข้อจำกัดที่ยังไม่ได้แก้ในโครงงานนี้


เหตุผลเชิงทฤษฎีที่สัญญาณนี้ได้ผลคือ\textbf{การนับคำซ้อนไม่ได้มาจากเวกเตอร์} จึงผิดพลาดคนละแบบกับสัญญาณเดิมและเสริมกันได้ ซึ่งสอดคล้องกับงานที่พบว่าการผสมสัญญาณแบบจับคู่คำเข้ากับการค้นด้วยเวกเตอร์ให้ผลดีกว่าใช้อย่างใดอย่างหนึ่ง \cite{karpukhin2020dpr, luan2021sparse} และสอดคล้องกับการประเมินข้ามโดเมนที่พบว่าแบบจำลองเวกเตอร์เสียเปรียบเมื่อคำค้นต้องการการจับคู่คำตรงตัว \cite{thakur2021beir} ข้อจำกัดคือฝั่งไทยไม่ได้ประโยชน์จากสัญญาณนี้ เพราะนับคำซ้อนกับคำแปล ไม่ใช่คำที่ผู้ใช้พิมพ์จริง ความคลาดเคลื่อนของชั้นแปลจึงไหลลงมาทำลายสัญญาณปลายน้ำ

\subsection{ความสามารถที่แท้จริงของการค้นซ้ำ}

เมื่อตรวจสอบว่าการค้นซ้ำช่วยอย่างไร พบว่า\textbf{ไม่ใช่การแก้คำสะกดผิดอย่างที่เข้าใจตอนแรก} การทดสอบกับคำค้นที่สะกดผิดพบว่ากู้คืนได้เพียง 3 จาก 12 คำ กรณีที่สำเร็จเป็นชื่อแบรนด์ดังที่แบบจำลองรู้จักอยู่แล้ว

ความสามารถจริงของการค้นซ้ำคือ\textbf{การเปลี่ยนไปใช้คำที่คลังสินค้าใช้เรียก} ซึ่งเป็นงานทางภาษาที่แบบจำลองทำได้ดี ต่างจากการตัดสินความตรงซึ่งแบบจำลองขนาดนี้ทำไม่ได้

% -----------------------------------------------------------
\section{ผลเชิงลบที่รายงานไว้}
% -----------------------------------------------------------

โครงงานนี้ทดลองแนวทาง 3 อย่างที่คาดว่าจะช่วย แต่วัดแล้วไม่ได้ผล จึงรายงานไว้พร้อมหลักฐาน เนื่องจากแต่ละข้อตัดความเป็นไปได้ทางหนึ่งออกด้วยการวัด ไม่ใช่การคาดเดา

\subsection{ให้เอเจนต์ปรับน้ำหนักของสัญญาณเอง}

ทดลองให้แบบจำลองภาษาเพิ่มน้ำหนักของสัญญาณการนับคำซ้อนเมื่อเห็นว่าผลลัพธ์ยังไม่ตรง ผลคือ\textbf{การตั้งน้ำหนักเป็น 20 40 และ 100 ให้ลำดับผลลัพธ์เหมือนกันทุกประการ} ซึ่งเป็นหลักฐานโดยตรงว่าวิธีรวมอันดับที่ใช้นั้น\textbf{อิ่มตัว} เพราะใช้เฉพาะอันดับไม่ใช้ค่าดิบ นอกจากนี้ เมื่อคะแนนการนับคำซ้อนต่ำ ย่อมแปลว่าสินค้าในกองไม่มีคำที่ผู้ใช้ขอมาตั้งแต่แรก การจัดอันดับใหม่จึงสร้างสินค้าที่ไม่มีขึ้นมาไม่ได้

\subsection{ให้เอเจนต์ระบุว่าเงื่อนไขข้อใดที่ผลลัพธ์ยังไม่ตอบ}

ทดลองให้แบบจำลองภาษาอ่านคำค้นและผลลัพธ์ แล้วบอกว่าเงื่อนไขข้อใดยังไม่ถูกตอบ เพื่อนำคำนั้นไปถ่วงน้ำหนักให้หนักขึ้น ผลคือ\textbf{แย่ลง 1.1 ถึง 3.6 จุด} สาเหตุคือแบบจำลองทิ้งคำที่ควรเก็บ เช่น ทิ้งคำที่เป็นชื่อแบรนด์เพราะเข้าใจว่าไม่ใช่เงื่อนไข และเก็บคำที่ไม่ใช่เงื่อนไข เมื่อเทียบกันแล้วพบว่า\textbf{การหาคำที่ขาดด้วยกฎในโปรแกรมแม่นยำกว่าการให้แบบจำลองหา} โดยทำได้ถูกต้อง 73 ถึง 83 เปอร์เซ็นต์ของคำค้น

ผลข้อนี้เป็น\textbf{หลักฐานชิ้นที่สาม}ที่ชี้ไปทางเดียวกันว่าแบบจำลองขนาด 8 พันล้านพารามิเตอร์ตัดสินความตรงของสินค้าไม่ได้ ต่อจากการวัดความสอดคล้องของแบบจำลองในบทบาทผู้ตัดสินที่ได้ค่าเพียง 0.08 ถึง 0.25 และการให้แบบจำลองตรวจแก้งานตัวเองที่ล้มเหลว

\subsection{บังคับให้ค้นซ้ำเมื่อชื่อสินค้าไม่มีคำที่ผู้ใช้ขอ}

ทดลองเพิ่มกฎว่าหากชื่อสินค้าที่เลือกได้ไม่มีคำที่ผู้ใช้ขอ ให้บังคับค้นซ้ำ ผลคือกฎนี้\textbf{เป็นที่มาของการแพ้ทั้งหมด} โดยทำให้ชนะ 5 แพ้ 4 ขณะที่กฎเดิมซึ่งดูเพียงว่าได้สินค้าครบหรือไม่ ชนะ 12 แพ้ 0

สาเหตุคือคลังสินค้าใช้คำเรียกต่างจากที่ผู้ใช้ใช้ กระถางที่ปลูกกุหลาบได้ไม่จำเป็นต้องมีคำว่ากุหลาบอยู่ในชื่อสินค้า

\textbf{บทเรียนที่ใช้ได้กับทั้งโครงงาน} คือการนับคำซ้อน\textbf{ใช้ได้ที่ชั้นจัดอันดับ แต่ใช้ไม่ได้ที่สวิตช์เปิดการค้นซ้ำ} เพราะการจัดอันดับเป็นการเปรียบเทียบของในกองเดียวกัน ส่วนการตัดสินใจค้นซ้ำต้องใช้เกณฑ์สัมบูรณ์

\subsection{ผลของการถอดกฎดังกล่าวออกเป็นการแลกเปลี่ยน}

การถอดกฎบังคับค้นซ้ำออกไม่ได้ดีขึ้นทุกด้าน แต่เป็นการแลกเปลี่ยน จึงรายงานทั้งสองชุดดังตารางที่ \ref{tab:tradeoff} ไม่เลือกรายงานเฉพาะชุดที่ให้ผลดี

\begin{table}[H]
\centering
\caption{เปรียบเทียบผลก่อนและหลังถอดกฎบังคับค้นซ้ำ}
\label{tab:tradeoff}
\begin{tabular}{|l|c|c|}
\hline
\textbf{รายการ} & \textbf{ยังเปิดกฎ} & \textbf{ถอดกฎออก (ที่รายงาน)} \\
\hline
อังกฤษ & ชนะ 9 \textbf{แพ้ 3} $\cdot$ $p = 0.0239$ & ชนะ 8 \textbf{แพ้ 0} $\cdot$ $p = 0.0078$ \\ \hline
ไทย & ชนะ 8 \textbf{แพ้ 1} $\cdot$ $p = 0.0195$ & ชนะ 5 \textbf{แพ้ 0} $\cdot$ $p = 0.0625$ \\ \hline
ปฏิเสธถูก อังกฤษ $\cdot$ ไทย & 48/60 $\cdot$ 49/60 & \textbf{50/60} $\cdot$ \textbf{51/60} \\
\hline
\end{tabular}
\end{table}

ฝั่งอังกฤษดีขึ้นทุกด้าน แต่ฝั่งไทยหลุดนัยสำคัญเพราะการค้นซ้ำทำงานน้อยลงจึงชนะน้อยลงตามไปด้วย สรุปคือเป็นการแลกระหว่างการชนะหลายข้อกับการไม่แพ้เลยแม้แต่ข้อเดียว โครงงานนี้เลือกอย่างหลังเพราะระบบที่ไม่เคยแย่ลงกว่าระบบพื้นฐานมีความเสี่ยงในการนำไปใช้จริงต่ำกว่า
